\subsection{Intersection properties}

LEMMA 7. Let E be a right A-module, F a left A-module and \(\mathbf{F}'\) , \(\mathbf{F}''\) two submodules of F such that \(\mathbf{F} = \mathbf{F}' + \mathbf{F}''\) . Then the intersection of the canonical images of \(\mathbf{E} \otimes \mathbf{F}'\) and \(\mathbf{E} \otimes \mathbf{F}''\) in \(\mathbf{E} \otimes \mathbf{F}\) is equal to the canonical image of \(\mathbf{E} \otimes (\mathbf{F}' \cap \mathbf{F}'')\) .

Consider the diagram

\[
\begin{array}{c} 0 \longrightarrow \mathrm{F} ^ {\prime} \cap \mathrm{F} ^ {\prime \prime} \longrightarrow \mathrm{F} ^ {\prime} \longrightarrow \mathrm{F} ^ {\prime} / (\mathrm{F} ^ {\prime} \cap \mathrm{F} ^ {\prime \prime}) \longrightarrow 0 \\ \Biggl \downarrow \quad \Biggl \downarrow \quad \text {   j   } \\ 0 \longrightarrow \mathrm{F} ^ {\prime \prime} \longrightarrow \mathrm{F} ^ {\prime} + \mathrm{F} ^ {\prime \prime} \longrightarrow (\mathrm{F} ^ {\prime} + \mathrm{F} ^ {\prime \prime}) / \mathrm{F} ^ {\prime \prime} \longrightarrow 0 \end{array}
\]

where the unspecified arrows are the canonical injections and surjections and j is the canonical isomorphism defined in Algebra, Chapter I, § 6, no. 13, Theorem 6. This diagram is commutative and its rows are exact. We derive (since \(F = F' + F''\) ) a commutative diagram

\[
\begin{array}{c} \mathbf {E} \otimes (\mathbf {F} ^ {\prime} \cap \mathbf {F} ^ {\prime \prime}) \longrightarrow \mathbf {E} \otimes \mathbf {F} ^ {\prime} \longrightarrow \mathbf {E} \otimes (\mathbf {F} ^ {\prime} / (\mathbf {F} ^ {\prime} \cap \mathbf {F} ^ {\prime \prime})) \\ \Biggl \downarrow \qquad \qquad \qquad \qquad \qquad \qquad \qquad \qquad \qquad \qquad \qquad \qquad \qquad 1 _ {\mathbf {E}} \otimes \mathbf {j} \\ \mathbf {E} \otimes \mathbf {F} ^ {\prime \prime} \longrightarrow \mathbf {E} \otimes \mathbf {F} \longrightarrow \mathbf {E} \otimes (\mathbf {F} / \mathbf {F} ^ {\prime \prime}) \end{array}
\]

The rows of this diagram are exact (no. 1, Lemma 1) and \(1_{E} \otimes j\) is an isomorphism. Our assertion is then a special case of § 1, no. 4, Proposition 1, (i) (cf.~Exercise 5).

PROPOSITION 6. Let E be a right A-module and F a left A-module such that E is F-flat. For every submodule \(\mathbf{F}'\) of F, we denote by \(\phi(\mathbf{F}')\) the image of \(\mathbf{E} \otimes \mathbf{F}'\) under the canonical mapping of \(\mathbf{E} \otimes \mathbf{F}'\) to \(\mathbf{E} \otimes \mathbf{F}\) (which is injective by Definition 1 of no. 2). Then, if \(\mathbf{F}'\) , \(\mathbf{F}''\) are two submodules of F, we have

\[
\phi (\mathbf {F} ^ {\prime} \cap \mathbf {F} ^ {\prime \prime}) = \phi (\mathbf {F} ^ {\prime}) \cap \phi (\mathbf {F} ^ {\prime \prime}).
\]

As E is F-flat, \(\phi(F' + F'')\) is identified with \(E \otimes (F' + F'')\) , and the submodules \(\phi(F')\) , \(\phi(F'')\) and \(\phi(F' \cap F'')\) are identified with the canonical images of \(E \otimes F'\) , \(E \otimes F''\) and \(E \otimes (F' \cap F'')\) in \(E \otimes (F' + F'')\) respectively. Proposition 6 then follows from Lemma 7.

Remark 1. With the hypotheses of Proposition 6, \(\mathbf{E} \otimes \mathbf{F}'\) is usually identified with \(\phi(\mathbf{F}')\) for every submodule \(\mathbf{F}'\) of \(\mathbf{F}\) , which gives the formula

\[
\mathrm{E} \otimes_ {\mathrm{A}} (\mathrm{F} ^ {\prime} \cap \mathrm{F} ^ {\prime \prime}) = (\mathrm{E} \otimes_ {\mathrm{A}} \mathrm{F} ^ {\prime}) \cap (\mathrm{E} \otimes_ {\mathrm{A}} \mathrm{F} ^ {\prime \prime}).
\]

PROPOSITION 7. Let E be a right A-module, E' a submodule of E, F a left A-module and F' a submodule of F. Suppose that E/E' or F/F' is flat. Then the canonical image of \(E' \otimes F'\) in \(E \otimes F\) is the intersection of the canonical images of \(E' \otimes F\) and \(E \otimes F'\) in \(E \otimes F\).

Suppose for example that \(E/E'\) is flat and consider the diagram

\[
\begin{array}{c} \mathrm{E} ^ {\prime} \otimes \mathrm{F} ^ {\prime} \longrightarrow \mathrm{E} \otimes \mathrm{F} ^ {\prime} \longrightarrow (\mathrm{E} / \mathrm{E} ^ {\prime}) \otimes \mathrm{F} ^ {\prime} \\ \Biggl \downarrow \qquad \qquad \qquad \qquad \qquad \qquad \qquad \qquad \qquad \qquad \qquad \qquad \qquad u \Biggl \downarrow \\ \mathrm{E} ^ {\prime} \otimes \mathrm{F} \longrightarrow \mathrm{E} \otimes \mathrm{F} \longrightarrow (\mathrm{E} / \mathrm{E} ^ {\prime}) \otimes \mathrm{F} \end{array}
\]

where the arrows are the canonical homomorphisms. This diagram is commutative and its rows are exact (no. 1, Lemma 1). As E/E' is flat, u is injective. Then our assertion is a special case of § 1, no. 4, Proposition 1, (i).

COROLLARY. Let E be a right A-module and E' a submodule of E.

\begin{enumerate}
\def\labelenumi{(\roman{enumi})}
\tightlist
\item
  Suppose that \(\mathbf{E} / \mathbf{E}'\) is flat. Then, for every left ideal \(\mathfrak{a}\) of \(\mathbf{A}\) ,
\end{enumerate}

\[
\mathbf {E} ^ {\prime} \mathfrak {a} = \mathbf {E} ^ {\prime} \cap \mathbf {E} \mathfrak {a}.\tag{5}
\]

\begin{enumerate}
\def\labelenumi{(\roman{enumi})}
\setcounter{enumi}{1}
\item
  Conversely, suppose that E is flat and, for every finitely generated left ideal a of A, relation (5) holds. Then E/E' is flat.
\item
  It is sufficient to apply Proposition to the case \(\mathbf{F} = \mathbf{A}_s\) , \(\mathbf{F}' = a\) .
\item
  To prove that \(E/E'\) is flat, apply criterion (c) of Proposition 4 of no. 5; it is then necessary to establish that the sequence
\end{enumerate}

\[
0 \rightarrow \mathrm{E} ^ {\prime} / \mathrm{E} ^ {\prime} a \rightarrow \mathrm{E} / \mathrm{E} a \rightarrow \mathrm{E} / (\mathrm{E} ^ {\prime} + \mathrm{E} a) \rightarrow 0
\]

is exact at E'/E'a for every finitely generated left ideal a of A. Now this is precisely what relation (5) expresses.

Remark 2. The conclusion of Proposition 7 remains true if we assume only that \(\mathbf{E} / \mathbf{E}'\) is F-flat or that \(\mathbf{F} / \mathbf{F}'\) is E-flat.

